\subsection{Graded regular rings}

Let \(A_{0}\) a ring and P a \(A_{0}\) graded module of type N with degrees \textgreater{} 0. Let A denote the ring. \(\mathsf{S}_{\mathrm{A}_{0}}(\mathrm{P})\) ; there exists on A a unique N-grading for which \(A_{0}\) is of degree 0 and P is a sub- \(A_{0}\) -graded module over A. Denote \(A_{+}\) the ideal \(\bigoplus_{n>0} A_{n}\) of A. Then the canonical map \(P \longrightarrow A_{+}/A_{+}^{2}\) is an isomorphism of \(A_{0}\) -graded modules (cf. A, III, p. 76, prop. 10).

If the \(A_{0}\) If the module P is graded free (A, II, p.~167, remark 3) and if \((x_{i})_{i\in I}\) is a basis of P consisting of homogeneous elements, the \(A_{0}\) -graded algebra \(\mathbf{S}_{\Lambda_{0}}(\mathbf{P})\) is isomorphic to the polynomial algebra \(\mathrm{A}_{0}[(\mathrm{X}_{i})_{i\in I}]\) equipped with the graduation for which each \(X_{i}\) is homogeneous of degree \(\deg(x_{i})\) We call \(A_{0}\) -graded polynomial algebra all \(A_{0}\) -graded algebra of type N, isomorphic to a \(A_{0}\) -graded algebra of the preceding form.

When the ring \(\Lambda_0\) is regular and the \(\mathrm{A}_0\) -module P is finitely generated projective, the ring \(\mathsf{S}_{\mathrm{A}_0}(\mathrm{P})\) is regular ( \(\S 4\) (see n° 2, cor. 4 of prop. 4). Conversely:

THEOREM 2.--- Let A be a regular ring, graded of type N. The ring A₀ formed by the elements of degree 0 in A is regular; there exists a finitely generated graded projective A₀-module P with degrees \textgreater{} 0 such that A is isomorphic as a graded A₀-algebra to \(S_{A_{0}}(P)\).

Let P denote the \(A_{0}\) graded module \(A_{+}/A_{+}^{2}\) . By prop. 4 of no. 3, the ring \(A_{0}\) is regular and the \(A_{0}\) -module P is projective and finitely generated. The homogeneous components of P are therefore projective, and there exists a section \(A_{0}\) -linear \(\varphi: P \to A_{+}\) , graded of degree 0, of the canonical surjection \(A_{+} \to P\) . Let \(f: S_{\mathrm{A}_{0}}(\mathrm{P}) \longrightarrow \Lambda\) the homomorphism of \(A_{0}\) -algebras which extends \(\varphi\) . By prop. 4, f extends to an isomorphism of the separated completion of \(\mathbf{S}_{\mathrm{A}_{0}}(\mathrm{P})\) for the topology \(\mathbf{S}_{\mathrm{A}_{0}}(\mathrm{P})_{+}\) -adically onto the separated completion of A for the topology \(A_{+}\) -adic. Consequently, f is injective and its image is dense in A for the topology \(A_{+}\) -adic. But since the induced topologies on the homogeneous components of A are discrete and the image of f is a graded submodule, this implies that f is bijective.

COROLLARY 1. - Let B be a regular ring, graded with positive degrees. Suppose that every \(\mathbf{B}_0\) -module projective of finite type is free.

\begin{enumerate}
\def\labelenumi{\alph{enumi})}
\item
  The ring B \(_{0}\) is integral and regular, and the B \(_{0}\) -algebra 13 is a B \(_{0}\) graduated polynomial algebra, of finite type.
\item
  Let A be a graded subring of B such that \(A_{0} = B_{0}\) and let B be a finitely generated A-module. The following conditions are equivalent:
  \begin{enumerate}
  \def\labelenumii{(\roman{enumii})}
  \item
  the A-module B is graded free;
  \item
  the A-module B is flat;
  \item
  A is a \(A_{0}\) graduated polynomial algebra, of finite type.
  \end{enumerate}
\end{enumerate}

By theorem 2, the ring B₀ is regular, hence a product of regular integral domains (§ 4, no. 2, example 2). For every idempotent element e of B₀, the B₀-module B₀e is projective, hence free, which implies e = 0 or 1; thus B₀ is an integral domain. Assertion a) then follows from theorem 2 and from the fact that a graded projective B₀-module of finite type is graded free.

Let us prove b).

\begin{enumerate}
\def\labelenumi{(\roman{enumi})}
\item
  \(\Leftrightarrow\) (ii): the \(A_0\) -graded flat modules of finite type are graded free (II, § 5, no. 2, cor. 2 of th. 1). The equivalence of (i) and (ii) therefore follows from A, X, p.~144, prop. 8.
\item
  \(\Leftrightarrow\) (iii): since B is integral over A and is a finitely generated A \(_0\) -algebra, A is an A \(_0\) -algebra of finite type (V, § 1, no. 9, lemma 5), hence a Noetherian ring. The equivalence of (ii) and (iii) then follows from the corollary of prop. 8 of § 4, no. 5 and from a) applied to A.
\end{enumerate}

COROLLARY 2.- Let \(k\) be a field, B a \(k\) -graded polynomial algebra, of finite type, and A a graded subalgebra of B. The following conditions are equivalent:

\begin{enumerate}
\def\labelenumi{(\roman{enumi})}
\item
  B is a graded free Λ-module;
\item
  B is a flat A-module;
\item
  ; we have \(\mathrm{Tor}_1^\Lambda (k,\mathbf{B}) = 0\) ;
\item
  the algebra A is a finitely generated graded polynomial k-algebra, and every algebraically free generating sequence of A consisting of homogeneous elements is B-regular.
\end{enumerate}

The implications (i) \(\Rightarrow\) (ii) and (ii) \(\Rightarrow\) (iii) are clear, and implication (iii) \(\Rightarrow\) (i) follows from \(\Lambda\) , X, p.~144, prop. 8, a).

(i)⇒(iv): since the ring B is regular and faithfully flat over A, the ring A is Noetherian by Prop. 11 of I, § 3, No. 6, and regular by Prop. 8 of § 4, No. 5, hence is a graded polynomial k-algebra (Cor. 1, a)). Every algebraically free generating sequence of A is A-regular (A, X, p. 158, example), hence B-regular since B is flat over A.

((iv) \(\Rightarrow\) (iii)): Suppose condition (iv) is satisfied, and let x be an algebraically free generating sequence of A consisting of homogeneous elements. Since the sequence x is A-regular, it is completely secant for A (A, X, p.~157, prop. 5), so that the A-module \(\mathrm{Tor}_{1}^{\mathrm{A}}(k,\mathrm{B})\) is isomorphic to \(\mathrm{H}_{1}(\mathbf{x},\mathrm{B})\) (A, X, p.~159, remark 3); but the latter is zero, since the sequence x is B-regular.

Corollaries 1 and 2 imply Lemma 5 of LIE, V, § 5, No. 5.

